\section{Collision divisors and elliptic tails}\label{p8-sec:hurwitz}

Fix a smooth compact curve \(C\) of genus \(g\ge 2\), put
\(H_m=1+m(g-1)\), and let \(m\ge 2\). Inside the stack of connected
stable maps \(\overline{\mathcal M}_{H_m+1}(C,m)\), let
\(\mathcal F_m=\mathcal F_m(C)\) be the locus of maps finite over \(C\).
These maps are flat, since the source is a nodal curve and its
structure sheaf is torsion-free over the smooth target. The locus is
open by properness of the universal curve and openness of
quasi-finiteness. Its complement is a union of open and closed
elliptic-tail loci, as proved in Section~\ref{p8-sec:tails}. Thus
\(\mathcal F_m\) is proper. Bryan and Pandharipande introduced this
surface and its branch morphism in their calculation of local
invariants of curves \cite[\S\S1.3.2, 4]{bp01}.

For a finite map, Riemann--Hurwitz on the normalization and the genus
formula for a nodal curve give
\[
           2=\deg R+2\,\#\{\text{nodes}\},
\]
where \(R\) is the ramification divisor on the normalization. Thus a
nodal map has one node and its normalization is \'etale. At a smooth
ramification point of index \(b\) the local presentation \(w=z^b\)
identifies the relative cotangent complex with the module
\(\mathcal O/(z^{b-1})\) in degree zero. At the node there are
coordinates \(uv=0\), \(w=u+v\), since both branches are unramified.
Writing its local algebra over the target as
\(\mathcal O_C[u]/(u^2-wu)\) gives the relative cotangent complex
\(\mathcal O/(2u-w)\) in degree zero. Here \(2u-w=u-v\) is a
nonzerodivisor and the quotient has length two. In each case the
module has projective dimension one. Its sheaf \(\mathcal{H}\!om\)
to \(\mathcal O\) vanishes, and its sheaf \(\mathcal{E}\!xt^1\) has the
same finite length as the module. The local-to-global sequence for
deformations over \(C\) consequently gives tangent dimension two and
zero obstruction. This proves smoothness of the stack and identifies
its virtual class with its fundamental class.

The branch morphism
\(\beta:\mathcal F_m\to C^{(2)}\) sends a smooth cover to its branch
divisor and a nodal one to twice the image of the node. Let
\(\Delta=\{2x:x\in C\}\). The connected components of the reduced
preimage of \(\Delta\) are of three kinds.
\begin{itemize}
\item \(N_i\) consists of nodal covers. The two points above the node
  belong to a possibly disconnected \'etale normalization.
\item \(T_j\) consists of smooth covers having a single ramification
  point of index three.
\item \(U_k\) consists of smooth covers having two distinct simple
  ramification points over the same point of \(C\).
\end{itemize}
For any such component \(B\), write \(a_B:B\to C\) for its branch
value and \(d_B=\int_B a_B^*[{\rm pt}]\). This is a stack degree.
For a fixed branch value it sums \(1/|\operatorname{Aut}_C(f)|\)
over the corresponding isomorphism classes. Indeed, the monodromy
data over the complement of the branch value are discrete and finite
in fixed degree, and move uniquely up to isomorphism in a disk of
branch values. In the nodal case the data consist of an \'etale cover
and an unordered pair of distinct sheets to be identified, with
connected quotient. Automorphisms interchanging the branches are
included.

In a versal deformation of a nodal map the local model is
\(uv=t,\ w=u+v+v_0\). Its two critical values have squared difference
\(16t\). At a point of \(T_j\) the model is \(w=z^3+uz+v\), for
which the squared difference is a nonzero constant times \(u^3\).
At a point of \(U_k\) the two critical values \(v_1,v_2\) are
independent parameters on an ordering cover; their squared
difference is \((v_1-v_2)^2\). The local deformation calculation
above gives all the parameters in these models. In particular the
reduced divisors are smooth and pairwise disjoint, and
\begin{equation}\label{p8-eq:hurwitz-pullback}
        \beta^*\Delta=\sum_i N_i+3\sum_j T_j+2\sum_k U_k.
\end{equation}
The multiplicities are orders on stack charts. Since the pullback of
\(\Delta\) to \(C\times C\) is twice the diagonal, \(\Delta^2=4-4g\).
Restricting \eqref{p8-eq:hurwitz-pullback} to a component of multiplicity
\(\ell_B\) gives
\begin{equation}\label{p8-eq:hurwitz-intersection}
                         B^2=\frac{4-4g}{\ell_B}\,d_B .
\end{equation}
These are the branch intersections used in \cite[\S4.3]{bp01}.

For a theta characteristic \(\theta\) define the following signs on the
reduced divisors. On \(N_i\), let \(\widetilde f\) be the normalization
of the cover. On \(T_j\), write \(p\) for the ramification point. Put
\begin{equation}\label{p8-eq:hurwitz-signs}
 s_\theta(N_i)=(-1)^{h^0(\widetilde f^*\theta)},\qquad
 s_\theta(T_j)=(-1)^{h^0(f^*\theta(p))}.
\end{equation}
In each expression the line is a theta characteristic on a smooth
curve, and its parity is locally constant. On a disconnected
normalization \(h^0\) is the sum over its components.

\begin{theorem}\label{p8-thm:boundary}
In \(H^2(\mathcal F_m;\QQ)\) the Euler class of the oriented spin
obstruction plane is
\begin{equation}\label{p8-eq:hurwitz-euler}
 e(V_\theta)=-\frac14
    \left(\sum_i s_\theta(N_i)[N_i]+\sum_j s_\theta(T_j)[T_j]\right),
\end{equation}
where brackets denote the classes of the reduced divisors.
For \(D=D_\theta\oplus D_\eta\), the finite-map contribution is
\begin{equation}\label{p8-eq:hurwitz-finite}
 F_m(D)=-\frac{g-1}{4}\left(
    \sum_i d_{N_i}s_\theta(N_i)s_\eta(N_i)
    +\frac13\sum_j d_{T_j}s_\theta(T_j)s_\eta(T_j)\right).
\end{equation}
\end{theorem}

\begin{proof}
The cokernels over finite nodal maps extend the oriented plane on the
smooth locus by the normal obstruction bundle construction
\cite[Theorem~8.1 and Lemma~11.1]{lpob}. Off the \(N_i\) and \(T_j\) there are two
distinct simple ramification points. By Theorem~\ref{p8-thm:residue}
their residue spaces are orthogonal lines. Ordering the points and
choosing a unit vector for each complex quadratic form gives a finite
cover trivializing the orthogonal bundle on this entire open set,
including the \(U_k\). Uniqueness up to homotopy of its positive
real form and transfer imply that the rational Euler class restricts
to zero there. By the Thom isomorphism for the disjoint smooth
divisors, it is therefore a linear combination of the classes
\([N_i]\) and \([T_j]\).

Set \(\kappa=c_1(K_C)\). On \(N_i\) the normalization sequence of
Corollary~\ref{p8-cor:nodal-spin-matching} gives
\(e(V_\theta)|_{N_i}=s_\theta(N_i)a_{N_i}^*\kappa/2\).
On \(T_j\) Theorem~\ref{p8-thm:profiles} gives
\(e(V_\theta)|_{T_j}=s_\theta(T_j)a_{T_j}^*\kappa/6\).
Integrating over a component \(B\) and dividing by its nonzero
self-intersection \eqref{p8-eq:hurwitz-intersection} determines its
coefficient. It is \(-s_\theta(B)/4\) in both cases. Finally the
oriented cokernel for \(D\) is the direct sum of the two spin
obstruction planes. Integrating the product of their Euler classes
and applying \eqref{p8-eq:hurwitz-intersection} gives
\eqref{p8-eq:hurwitz-finite}.
\end{proof}
\subsection{Contracted elliptic tails}\label{p8-sec:tails}

Let \(N\to C\) be a complex vector bundle of rank two and degree
\(2g-2\), and \(D\) a real Cauchy--Riemann operator on \(N\). Assume
that the pullback operator on every smooth connected cover of \(C\)
is injective. On an \'etale cover the index is zero, so the operator
is invertible. Orient all normal determinants by continuation from
complex-linear operators, and denote the determinant sign of an
\'etale cover \(f\) by \(\sigma_D(f)\). The signed stack degree of
the connected \'etale covers of degree \(m\) is
\[
 c_m(D)=\sum_{[f:Y\to C]}
                  \frac{\sigma_D(f)}{|\operatorname{Aut}_C(f)|}.
\]
Write \(L_{m,H_m+1}(D)\) for the absolute normal Euler number on
\(\overline{\mathcal M}_{H_m+1}(C,m)\) and \(F_m(D)\) for its
integral over the finite-map locus. Put \(F_1(D)=0\), since a
finite degree-one cover is an isomorphism. When \(D\) is the normal
operator of a super-rigid embedded curve, the fixed-target
comparison identifies this absolute Euler pairing with the local
Gromov--Witten degree \cite[Theorem~1.2(2)--(3)]{zinger}.
For $m\ge2$ the relatively oriented Euler number of
\cite[Definition~6.2 and Remark~6.4]{bai} differs from this absolute
number by the determinant sign of the operator on $C$.


\begin{proposition}\label{p8-prop:tails}
The complement of the finite-map locus is a disjoint union of
open and closed components indexed by the connected degree-\(m\)
\'etale covers \(f:Y\to C\). Each component has its ordinary
fundamental class as virtual class. Their contributions give
\begin{equation}\label{p8-eq:tails-identity}
         L_{m,H_m+1}(D)=F_m(D)-\frac{m(g-1)}{12}\,c_m(D).
\end{equation}
\end{proposition}

\begin{proof}
Consider a stable map, and let \(s\) be the number of nonconstant
components on the normalization of its domain. Write \(R\) for the
sum of their ramification divisors, \(g_0\) for the sum of the
genera of the normalized contracted components, and \(\Gamma\) for
the connected dual graph. Riemann--Hurwitz and the arithmetic genus
give
\[
             2=s+\frac{\deg R}{2}+g_0+b_1(\Gamma).
\]
If \(s=2\), both positive-degree components are \'etale and the
graph is a tree. A contracted rational vertex must have valence at
least three, which is impossible in a tree with at most two leaves.
If \(s=1\) and \(\deg R=2\), the same argument excludes a contracted
subtree. If \(s=1\) and \(\deg R=0\), the nonconstant component has
smooth normalization \(Y\) of genus \(H_m\), and \(g_0+b_1(\Gamma)=1\).
When \(g_0=1\) the tree consists of \(Y\) joined directly to a
smooth genus-one component. When \(b_1(\Gamma)=1\), every contracted
vertex has valence at least three. The identity
\(\sum_v(\operatorname{val}(v)-2)=0\) leaves just a self-node on
\(Y\), or one contracted rational vertex with a self-loop joined
to \(Y\) by one edge. Thus every map with a contracted component
consists of a smooth connected \'etale cover with a one-pointed
stable curve of genus one attached.

For fixed \(f\), gluing at \(x\in Y\) defines a proper morphism from
\[
 \mathcal G_f=
  \big[(Y\times\overline{\mathcal M}_{1,1})/\operatorname{Aut}_C(f)\big]
      \longrightarrow\overline{\mathcal M}_{H_m+1}(C,m).
\]
The deck group acts on \(Y\), and the pointed tail retains its stack
stabilizers. To compute the deformation theory at
\(Y\cup_{x=y}E\), first preserve the attaching node. The cover is
rigid because \(df:T_Y\to f^*T_C\) is an isomorphism. Its marked
point moves freely. The tail deforms as a pointed stable curve with
constant image in \(C\). Varying this image makes the matching
equation at the node transverse. The remaining obstruction is
\(T_{f(x)}C\otimes H^1(E,\mathcal O_E)\). The boundary map from
the smoothing line is
\begin{equation}\label{p8-eq:tails-smoothing}
 T_xY\otimes T_yE\ \xrightarrow{\,df_x\otimes\delta_y\,}\
 T_{f(x)}C\otimes H^1(E,\mathcal O_E),
\end{equation}
up to the choice of sign for the connecting homomorphism. Indeed,
in the plumbing \(uv=t\), the linear term on the cover gives the
principal part \(df_x(t/v)\) on the tail. Thus \(\delta_y\) is the
connecting homomorphism for
\(0\to\mathcal O_E\to\mathcal O_E(y)\to T_yE\to0\). The dualizing
sheaf is trivial for both a smooth elliptic curve and an irreducible
one-nodal genus-one curve. Serre duality and Riemann--Roch give
\(H^1(\mathcal O_E(y))=0\) and \(H^0(\mathcal O_E(y))=\CC\).
Hence \eqref{p8-eq:tails-smoothing} is an isomorphism on all of
\(\mathcal G_f\).

The full stable-map theory is therefore unobstructed of complex
dimension two along \(\mathcal G_f\), and its tangent space consists
of the attachment and pointed-tail deformations. These are precisely
the tangents to the source of the gluing morphism. Gluing identifies
the stabilizers and is injective on geometric points, since the
positive-degree component and the attaching node are intrinsic. It
is thus a proper representable \'etale open immersion, and its image
is closed as well. This proves the asserted component description
and the fundamental-class claim. It also gives the properness used in
Section~\ref{p8-sec:hurwitz}. An alternative description by lifting to
the maximal \'etale cover appears in \cite[\S2.2]{bp01}.

It remains to orient the normal obstruction. On \(Y\) the operator
\(f^*D\) is invertible. On the constant tail the normal operator
is \(\bar\partial\) with coefficients in \(N_{f(x)}\). Its kernel
consists of the constants, which evaluate isomorphically onto the
matching fiber. The normalization exact sequence at the attaching
node consequently identifies the full normal cokernel with
\[
                    H^1(E,\mathcal O_E)\otimes N_{f(x)}.
\]
This uses the arithmetic \(H^1\) also for a nodal tail. In families
the first factor is the dual of the Hodge line \(\mathbb E\). The
normalization isomorphism of determinants gives the constants and the
matching fiber their complex orientations; their determinants cancel.
Continuing from the complex-linear normalization sequence leaves
the factor \(\sigma_D(f)\) from the invertible cover operator
\cite[\S2]{zinger}, \cite[Appendix~B]{bai}. Thus the full normal
orientation is the complex orientation of
\(\mathbb E^\vee\otimes f^*N\) multiplied by \(\sigma_D(f)\).

Write \(\lambda=c_1(\mathbb E)\). On the product of the two curves
\[
        c_2(\mathbb E^\vee\otimes f^*N)=-\lambda\,c_1(f^*N).
\]
The genus-one discriminant has a simple zero at the nodal boundary
and weight twelve, so \(12\lambda\) is the boundary class on
\(\overline{\mathcal M}_{1,1}\). The boundary point has stabilizer
of order two, whence \(\int\lambda=1/24\). Finally
\(\int_Y c_1(f^*N)=m(2g-2)\). Dividing by
\(|\operatorname{Aut}_C(f)|\) and summing proves
\eqref{p8-eq:tails-identity}.
\end{proof}

We record the implication for the exponents of the elementary
factors. Let \(n_{d,a}\) be the coefficients in the expansion of
the connected local series in the connected elementary series
\(G_a(q^d,t)\), with \(a\ge H_d\). Here $G_a$ is the
connected Gromov--Witten series of the elementary genus-$a$ curve,
the logarithm of its normalized partition function. A connected
degree-$j$, genus-$h$ contribution has monomial $q^jt^{2h-2}$.
The same counts are attached to the balanced universal factors
in Section~\ref{p8-sec:exponents}. Write \(C_j(a)\) for the
unsigned connected \'etale degree-\(j\) count of a genus-\(a\)
curve, and \(F_j^{\rm elem}(a)\) for the finite-map Euler number of
its doubled spin operator at first excess. All \'etale signs of this
operator are positive. The elementary series begins in degree \(j\)
at genus \(1+j(a-1)\), and its degree-one leading coefficient is one
\cite[Propositions~3.2--3.3]{ip}, \cite[Proposition~3.46]{diw}. At genus
\(H_m\), comparison therefore gives
\[
               c_m(D)=\sum_{d\mid m} n_{d,H_d} C_{m/d}(H_d).
\]
At genus \(H_m+1\), a proper divisor \(d\mid m\) can still contribute
only the factor of genus \(H_d\), since increasing that genus by one
raises the least source genus by \(m/d\). For the equal-degree
factor of genus \(H_m\) the finite term is \(F_1^{\rm elem}(H_m)=0\).
In each divisor term the tail multiplier in
\eqref{p8-eq:tails-identity} is the same, since
\((m/d)(H_d-1)=m(g-1)\). The least-genus identity cancels their sum
against the tail contribution on the left and yields
\begin{equation}\label{p8-eq:tails-exponent}
          n_{m,H_m+1}=F_m(D)-
          \sum_{\substack{d\mid m\\d<m}}
                   n_{d,H_d}\,F_{m/d}^{\rm elem}(H_d).
\end{equation}
